\subsection{Construction-error recurrence after recorded corrective requests}
\label{sec:recurrence}

We extend the within-decision recovery analysis to later construction submissions in the same player-game. The population contains 422 decisions with a parsed construction call, including the 114 decisions with an uneven-intermediate-operation error. For each error, we identify the next distinct decision by that player containing a construction submission. No link crosses a game boundary. This conditions on submitting another build, not on receiving an identical investment opportunity.

In fifty-four error decisions, a later same-decision request contains the exact validator error in its delivered scenario notes. Fifty have a subsequent construction-submission decision; nine of those repeat the error on the first construction submission. Eighteen of the fifty next submissions remain within the same global turn. Among the thirty-two occurring on a later turn, eight repeat the error. Restricting to the first corrective-request episode per player-game yields thirty-one episodes, twenty-nine with a later submission, and two repetitions. These are overlapping observational sequences under changing states and allocations. The restricted count reduces repeated episode weighting, but does not create a randomized feedback contrast.

\begin{table}[t]
\centering\small
\begin{tabularx}{\linewidth}{Yrrr}
\toprule
Episode population & Episodes & Later submission & Error repeated\\
\midrule
All uneven-error decisions & 114 & 109 & 63\\
Exact error in a later corrective request & 54 & 50 & 9\\
First such request per player-game & 31 & 29 & 2\\
\bottomrule
\end{tabularx}
\caption{Construction-error recurrence at the next distinct submitted-build decision. Repetition concerns the first construction submission there. Missing later submissions are censored, not successful repairs. Populations overlap; the differences cannot identify a feedback effect.}
\label{tab:recurrence}
\end{table}

The configurations have very different access to a corrective request. Sixty Opus error decisions follow a refusal that already consumed an attempt; there is no later request carrying the construction error. The exceptional Opus decision, G2 d129, has the error before the refusal, so its later request contains feedback but the returned response refuses. A request recorded as sent is not proof that the provider processed its content as an agent decision. Among all error episodes with a later submitted build, fifty-four of the sixty-three repetitions belong to Opus. Describing that aggregate as an inability to learn from equivalent feedback would erase the treatment difference.

Receiving a corrective request and repairing the current action are also different conditions. Forty-eight of the fifty-four feedback episodes end in an applied build. Here repair means that some construction is applied; it does not require preserving the original purchase multiset. Forty-four have a next submission, and four repeat the error. The six episodes without an applied build all have a next submission, with five repetitions. Thus five of the nine repetitions in the request-exposure population follow an unrepaired decision. These nested descriptions answer a bounded adaptation question: some configurations with a recorded corrective request repeat a specific encoding error, including some that previously ended in an applied build. They do not show durable learning, improvement relative to no feedback, or transfer to unfamiliar opponents. Immediate repair, later error recurrence, retention of a rule, and improvement in strategy are distinct endpoints. The retained episode ledger records decision and prompt references, turn gaps, response conditions, censoring and actual applied actions.
